\documentclass[12pt,a4paper]{article}
\usepackage[utf8]{inputenc}
\usepackage[polish]{babel}
\usepackage[T1]{fontenc}
\usepackage{graphicx}
\usepackage{geometry}
\usepackage{listings}
\usepackage{xcolor}
\usepackage{float}
\usepackage{tikz}
\usepackage{hyperref}
\usepackage{enumitem}
\usepackage{amsmath}
\usepackage{tcolorbox}

\usetikzlibrary{shapes.geometric, shapes.multipart, arrows, positioning, fit, backgrounds}

\geometry{margin=2.5cm}

\definecolor{codegreen}{rgb}{0,0.6,0}
\definecolor{codegray}{rgb}{0.5,0.5,0.5}
\definecolor{codepurple}{rgb}{0.58,0,0.82}
\definecolor{backcolour}{rgb}{0.95,0.95,0.92}

\lstdefinestyle{pythonstyle}{
    backgroundcolor=\color{backcolour},   
    commentstyle=\color{codegreen},
    keywordstyle=\color{magenta},
    numberstyle=\tiny\color{codegray},
    stringstyle=\color{codepurple},
    basicstyle=\ttfamily\footnotesize,
    breakatwhitespace=false,         
    breaklines=true,                 
    captionpos=b,                    
    keepspaces=true,                 
    numbers=left,                    
    numbersep=5pt,                  
    showspaces=false,                
    showstringspaces=false,
    showtabs=false,                  
    tabsize=2,
    literate={ą}{{\k{a}}}1 {ć}{{\'c}}1 {ę}{{\k{e}}}1 {ł}{{\l{}}}1 {ń}{{\'n}}1 {ó}{{\'o}}1 {ś}{{\'s}}1 {ź}{{\'z}}1 {ż}{{\.z}}1
             {Ą}{{\k{A}}}1 {Ć}{{\'C}}1 {Ę}{{\k{E}}}1 {Ł}{{\L{}}}1 {Ń}{{\'N}}1 {Ó}{{\'O}}1 {Ś}{{\'S}}1 {Ź}{{\'Z}}1 {Ż}{{\.Z}}1
}

\lstset{style=pythonstyle}

\title{\textbf{Dokumentacja Architektury Systemu IO\_RNG}\\
\large System Testowania i Porównywania Generatorów Liczb Pseudolosowych}
\author{Projekt IO}
\date{\today}

\begin{document}

\maketitle
\tableofcontents
\newpage

\section{O projekcie}

\subsection{Cel i przeznaczenie aplikacji}

IO\_RNG to zaawansowana aplikacja webowa dedykowana profesjonalnemu testowaniu i kompleksowemu porównywaniu różnorodnych algorytmów generowania liczb pseudolosowych (PRNG - Pseudo Random Number Generators), w tym także algorytmów kryptograficznie bezpiecznych (CSPRNG - Cryptographically Secure Pseudo Random Number Generators).

\subsubsection{Motywacja}

Generatory liczb pseudolosowych stanowią fundamentalny element wielu dziedzin informatyki i matematyki stosowanej. Ich jakość bezpośrednio wpływa na bezpieczeństwo systemów kryptograficznych, rzetelność symulacji Monte Carlo, wiarygodność algorytmów randomizowanych oraz poprawność metod statystycznych. Mimo powszechnego zastosowania, brakuje narzędzi umożliwiających łatwe i kompleksowe porównanie różnych implementacji algorytmów PRNG pod kątem zarówno jakości generowanych liczb, jak i wydajności obliczeniowej.

Aplikacja IO\_RNG odpowiada na tę potrzebę, oferując zunifikowane środowisko do testowania algorytmów PRNG zaimplementowanych w różnych językach programowania, wykonywania na nich standaryzowanych testów statystycznych oraz wizualizacji wyników w formie interaktywnych wykresów i profesjonalnych raportów PDF.

\subsubsection{Zakres funkcjonalny}

Aplikacja koncentruje się na trzech głównych aspektach:

\begin{itemize}
    \item \textbf{Generowanie} - produkcja sekwencji bitów lub liczb z zadanego zakresu z wykorzystaniem wybranego algorytmu PRNG wraz z konfigurowalnymi parametrami
    \item \textbf{Testowanie} - wykonywanie baterii testów statystycznych (NIST, Diehard, testy częstości i jednostajności) na wygenerowanych danych lub ciągach dostarczonych przez użytkownika
    \item \textbf{Analiza} - wizualizacja wyników, porównywanie algorytmów pod względem jakości i wydajności, generowanie raportów z pełną dokumentacją testów
\end{itemize}

\subsection{Przewodnik użytkownika}

Aplikacja została zaprojektowana z myślą o intuicyjności i elastyczności. Poniżej przedstawiono funkcjonalność poszczególnych komponentów systemu.

\subsubsection{Generator}

Moduł generatora stanowi centralny punkt interakcji z algorytmami PRNG. Podzielony jest na dwa tryby pracy:

\paragraph{Generator bitów}
Umożliwia wygenerowanie określonej liczby bitów (0 lub 1) przy użyciu wybranego algorytmu. Użytkownik może:
\begin{itemize}
    \item Wybrać algorytm PRNG z dostępnej listy (LCG, Park-Miller, PCG32, Mersenne Twister, Xoshiro256, ChaCha20, Blum-Blum-Shub, etc.)
    \item Określić liczbę bitów do wygenerowania
    \item Podać parametry startowe algorytmu (seed, parametry specyficzne dla danego generatora)
    \item Skonfigurować zaawansowane parametry:
    \begin{itemize}
        \item \texttt{bits\_per\_value} - liczba bitów wyodrębnianych z każdej wygenerowanej wartości
        \item \texttt{msb\_first} - kierunek ekstrakcji bitów (od najbardziej lub najmniej znaczącego bitu)
    \end{itemize}
\end{itemize}

Po wygenerowaniu, aplikacja prezentuje wizualizację strumienia bitów wraz z czasem wykonania operacji.

\paragraph{Generator liczb}
Pozwala na generowanie liczb całkowitych z wybranego zakresu. Dodatkowo oferuje:
\begin{itemize}
    \item Interaktywny wykres słupkowy dystrybucji wygenerowanych liczb
    \item Obliczanie i wyświetlanie średniej arytmetycznej
    \item Opcjonalną animację procesu generowania z efektami wizualnymi
    \item Wszystkie funkcje konfiguracji dostępne w generatorze bitów
\end{itemize}

\subsubsection{Testowanie algorytmów}

Moduł testów umożliwia przeprowadzenie kompleksowej analizy statystycznej algorytmów PRNG. System oferuje dwa źródła danych wejściowych:

\begin{itemize}
    \item \textbf{Testowanie algorytmu z listy} - wybór konkretnego generatora i automatyczne wygenerowanie wymaganej liczby próbek
    \item \textbf{Testowanie własnych bitów} - możliwość wklejenia sekwencji bitów dostarczonej przez użytkownika
\end{itemize}

Dostępne pakiety testów:
\begin{itemize}
    \item \textbf{Frequency Test} - test częstości występowania zer i jedynek
    \item \textbf{Uniformity Test} - test jednostajności rozkładu
    \item \textbf{NIST Test Suite} - pełny zestaw 15 testów opracowanych przez National Institute of Standards and Technology
    \item \textbf{Diehard Test Suite} - klasyczny zestaw testów George'a Marsgalii
\end{itemize}

Po wyborze testów i liczby próbek, system wyświetla pasek postępu wykonywanych testów w czasie rzeczywistym. Po zakończeniu prezentowane są:
\begin{itemize}
    \item Całkowity czas wykonania baterii testów
    \item Liczba testów zaliczonych i niezaliczonych
    \item Szczegółowa lista wyników z czasem wykonania każdego testu
    \item Wizualizacja wyników w formie wykresów
\end{itemize}

\subsubsection{Historia wyników}

Komponent historii przechowuje i zarządza wszystkimi wykonanymi testami. Funkcjonalności:
\begin{itemize}
    \item Lista wszystkich przeprowadzonych testów z podstawowymi informacjami (data, algorytm, typ testu, wynik)
    \item Podgląd szczegółów konkretnego testu (pełne statystyki, parametry użyte podczas testowania)
    \item Zaawansowane filtrowanie:
    \begin{itemize}
        \item Po algorytmie PRNG
        \item Po typie testu (NIST, Diehard, pojedyncze testy)
        \item Po dacie wykonania
        \item Po wyniku (zaliczone/niezaliczone)
    \end{itemize}
    \item Usuwanie pojedynczych wyników lub czyszczenie całej historii
\end{itemize}

\subsubsection{Dashboard i analiza}

Strona główna (Dashboard) stanowi centrum analityczne aplikacji, prezentujące zagregowane dane i umożliwiające głęboką analizę porównawczą:

\paragraph{Statystyki ogólne}
\begin{itemize}
    \item Liczba przetestowanych algorytmów
    \item Łączna liczba przeprowadzonych testów
    \item Interaktywne wykresy podsumowujące
\end{itemize}

\paragraph{Analiza wydajności}
Moduł pozwala na porównanie czasów wykonania różnych algorytmów. Dostępne są:
\begin{itemize}
    \item Wykresy porównawcze czasów generowania dla różnych rozmiarów próbek
    \item Szczegółowa analiza z dodatkowymi metrykami wydajnościowymi
    \item Wykres punktowy z możliwością dynamicznego odczytu parametrów
    \item Po wybraniu konkretnego algorytmu i testu - automatyczna transformacja do wykresu liniowego dla lepszej czytelności trendu
\end{itemize}

\paragraph{Analiza jakości testów}
Umożliwia ocenę jakości algorytmów pod kątem testów statystycznych:
\begin{itemize}
    \item Wizualizacja wskaźników zaliczalności poszczególnych testów
    \item Porównanie wyników różnych algorytmów na tym samym teście
    \item Identyfikacja słabych punktów algorytmów
\end{itemize}

\paragraph{Generator raportów PDF}
System oferuje zaawansowany mechanizm generowania profesjonalnych raportów w formacie PDF, zawierających:
\begin{itemize}
    \item Zestawienie wszystkich przeprowadzonych testów
    \item Wykresy porównawcze wydajności algorytmów
    \item Wykresy porównawcze jakości (wskaźniki zaliczalności testów)
    \item Szczegółowe statystyki dla każdego algorytmu
    \item Spersonalizowany motyw graficzny (zgodny z ustawieniami kolorystycznymi aplikacji)
\end{itemize}

\subsubsection{Wiki - baza wiedzy}

Aplikacja zawiera wbudowaną encyklopedię wiedzy o algorytmach PRNG i metodologii testowania:

\paragraph{Algorithms - Dokumentacja algorytmów}
Każdy dostępny w systemie algorytm PRNG posiada dedykowaną stronę dokumentacji zawierającą:
\begin{itemize}
    \item Podstawowe informacje i definicję algorytmu
    \item Matematyczny opis działania ze wzorami
    \item Kluczowe fragmenty implementacji w różnych językach programowania
    \item Kontekst historyczny i autorstwo
    \item Analiza złożoności obliczeniowej (czasowej i pamięciowej)
    \item Praktyczne przykłady obliczeniowe
    \item Zastosowania i ograniczenia
    \item Rekomendacje dotyczące parametrów
\end{itemize}

\paragraph{Methodology - Dokumentacja testów}
Szczegółowe opisy wszystkich dostępnych testów statystycznych:
\begin{itemize}
    \item Zasada działania testu
    \item Wzory matematyczne i podstawy teoretyczne
    \item Szczegóły implementacji w systemie
    \item Interpretacja wyników (wartości p-value, poziomy istotności)
    \item Parametry konfiguracyjne testu
    \item Bibliografia i źródła
\end{itemize}

\subsubsection{Ustawienia i personalizacja}

Moduł ustawień demonstruje elastyczność aplikacji pod względem dostosowania interfejsu użytkownika:

\begin{itemize}
    \item \textbf{Kolor wiodący} - możliwość wyboru głównego koloru motywu aplikacji
    \item \textbf{Kolor akcentów} - definiowanie koloru elementów interfejsu wyróżniających ważne funkcje
    \item \textbf{Tryb jasny/ciemny} - automatyczne dostosowanie kontrastu
    \item \textbf{Synchronizacja motywu} - wybrany schemat kolorystyczny jest automatycznie stosowany również w generowanych raportach PDF, zapewniając spójność wizualną
\end{itemize}

\subsubsection{O aplikacji}

Sekcja informacyjna zawierająca:
\begin{itemize}
    \item Podstawowe informacje o projekcie IO\_RNG
    \item Możliwość pobrania pełnej dokumentacji technicznej (niniejszy dokument)
    \item Informacje o twórcach i kontakt
    \item Licencja i warunki użytkowania
\end{itemize}

\subsection{Rozszerzenia względem założeń projektowych}

Aplikacja IO\_RNG przekroczyła pierwotne założenia projektowe, oferując szereg dodatkowych funkcjonalności:

\begin{itemize}
    \item Generator liczb z konfigurowalnymi zakresami (oprócz pierwotnie planowanego generatora bitów)
    \item Możliwość testowania własnych ciągów bitów dostarczonych przez użytkownika
    \item Zaawansowany system generowania raportów PDF z wykresami i metrykami
    \item Kompleksowa baza wiedzy Wiki z dokumentacją algorytmów i testów
    \item Przyjazny webowy interfejs użytkownika zaprojektowany zgodnie z zasadami UX/UI
    \item System personalizacji wyglądu z synchronizacją motywu między aplikacją a raportami
    \item Interaktywne wykresy z możliwością dynamicznego filtrowania i transformacji
    \item Animacje i efekty wizualne poprawiające doświadczenie użytkownika
\end{itemize}

\section{Wprowadzenie techniczne}

\subsection{Cel dokumentu}
Dokument przedstawia architekturę techniczną systemu IO\_RNG, zastosowane wzorce projektowe, zasady programowania obiektowego oraz wykorzystane struktury danych. Przeznaczony jest dla deweloperów i architektów oprogramowania zainteresowanych szczegółami implementacyjnymi systemu.

\subsection{Technologie}
\begin{itemize}
    \item \textbf{Backend:} Django REST Framework, Python 3.x
    \item \textbf{Frontend:} React, TypeScript, Vite
    \item \textbf{Baza danych:} SQLite
    \item \textbf{Generatory:} Python, C++, Rust, C\#
\end{itemize}

\section{Architektura Systemu}

\subsection{Przegląd architektury}

System IO\_RNG implementuje \textbf{Clean Architecture} (Czystą Architekturę) z wyraźnym podziałem na warstwy:

\begin{enumerate}
    \item \textbf{Warstwa domeny (Core)} - encje i logika biznesowa
    \item \textbf{Warstwa infrastruktury (Infrastructure)} - implementacje techniczne
    \item \textbf{Warstwa API} - interfejs REST
    \item \textbf{Warstwa prezentacji (Frontend)} - interfejs użytkownika
\end{enumerate}

\subsection{Schemat struktury systemu}

\begin{figure}[H]
\centering
\begin{tikzpicture}[
    node distance=0.8cm and 1.5cm,
    layer/.style={rectangle, draw=black!80, fill=blue!10, thick, minimum width=3cm, minimum height=1.2cm, align=center, rounded corners},
    component/.style={rectangle, draw=black!60, fill=green!10, thick, minimum width=2.8cm, minimum height=0.8cm, align=center},
    interface/.style={rectangle, draw=black!60, fill=yellow!20, thick, minimum width=2.8cm, minimum height=0.8cm, align=center},
    entity/.style={rectangle, draw=black!60, fill=red!10, thick, minimum width=2.8cm, minimum height=0.8cm, align=center},
    arrow/.style={thick,->,>=stealth}
]

% Frontend Layer
\node[layer] (frontend) at (0,0) {\textbf{Frontend Layer}\\React + TypeScript};

% API Layer
\node[layer, below=of frontend] (api) {\textbf{API Layer}\\Django REST Framework};
\node[component, below=0.3cm of api] (viewsets) {ViewSets};
\node[component, left=0.2cm of viewsets] (serializers) {Serializers};

% Core Layer
\node[layer, below=2cm of api] (core) {\textbf{Core Layer}\\Domain Logic};
\node[entity, below=0.3cm of core, xshift=-3cm] (entities) {Entities};
\node[interface, below=0.3cm of core] (interfaces) {Interfaces};
\node[component, below=0.3cm of core, xshift=3cm] (usecases) {Use Cases};

% Infrastructure Layer
\node[layer, below=2.5cm of core] (infra) {\textbf{Infrastructure Layer}\\Implementations};
\node[component, below=0.3cm of infra, xshift=-3cm] (models) {Django Models};
\node[component, below=0.3cm of infra] (repos) {Repositories};
\node[component, below=0.3cm of infra, xshift=3cm] (runners) {RNG Runners};

% Database
\node[layer, below=1.5cm of infra] (db) {\textbf{Database}\\SQLite};

% External Systems
\node[layer, right=3cm of infra] (external) {\textbf{Algorytmy RNG}\\Python/C++/Rust};

% Arrows
\draw[arrow] (frontend) -- (api);
\draw[arrow] (api) -- (core);
\draw[arrow] (core) -- (infra);
\draw[arrow] (infra) -- (db);
\draw[arrow] (runners) -- (external);
\draw[arrow, dashed] (interfaces) -- (repos);
\draw[arrow, dashed] (interfaces) -- (runners);

\end{tikzpicture}
\caption{Architektura warstwowa systemu IO\_RNG}
\end{figure}

\newpage

\subsection{Diagram klas - Core Layer}

\begin{figure}[H]
\centering
\begin{tikzpicture}[
    class/.style={rectangle split, rectangle split parts=3, draw=black!80, fill=blue!10, thick, minimum width=4cm, align=center},
    interface/.style={rectangle split, rectangle split parts=3, draw=black!80, fill=yellow!20, thick, minimum width=4cm, align=center},
    arrow/.style={thick,->,>=stealth},
    implements/.style={thick,dashed,->,>=stealth}
]

% Entities
\node[class] (rng) at (0,0) {
    \textbf{RNG (Entity)}
    \nodepart{second}
    \small
    - name: str\\
    - language: Language\\
    - algorithm: Algorithm\\
    - code path: str\\
    - parameters: Dict\\
    - is active: bool
    \nodepart{third}
    + validate for execution()\\
    + get parameter(key)
};

\node[class, right=2cm of rng] (testresult) {
    \textbf{TestResult (Entity)}
    \nodepart{second}
    \small
    - rng id: int\\
    - test name: str\\
    - passed: bool\\
    - score: float\\
    - execution time ms: float\\
    - statistics: Dict
    \nodepart{third}
    + post init()
};

% Enums
\node[class, below=2cm of rng] (language) {
    \textbf{Language (Enum)}
    \nodepart{second}
    \small
    PYTHON\\
    CPP\\
    RUST\\
    JAVA\\
    EXECUTABLE
    \nodepart{third}
};

\node[class, right=1cm of language] (algorithm) {
    \textbf{Algorithm (Enum)}
    \nodepart{second}
    \small
    LINEAR CONGRUENTIAL\\
    MERSENNE TWISTER\\
    XORSHIFT\\
    PCG\\
    SYSTEM\\
    CUSTOM
    \nodepart{third}
};

% Interfaces
\node[interface, below=5cm of rng, xshift=-2cm] (irngrepo) {
    \textbf{<<Interface>>}\\
    IRNGRepository
    \nodepart{second}
    \nodepart{third}
    \small
    + save(rng)\\
    + get by id(id)\\
    + get all()\\
    + delete(id)\\
    + update(rng)
};

\node[interface, right=1cm of irngrepo] (iresultrepo) {
    \textbf{<<Interface>>}\\
    ITestResultRepository
    \nodepart{second}
    \nodepart{third}
    \small
    + save(result)\\
    + get by id(id)\\
    + get by rng(rng id)\\
    + get all()\\
    + delete(id)
};

\node[interface, right=1cm of iresultrepo] (irunner) {
    \textbf{<<Interface>>}\\
    IRNGRunner
    \nodepart{second}
    \nodepart{third}
    \small
    + can run(rng)\\
    + generate numbers()\\
    + generate raw()\\
    + validate setup()
};

% Relationships
\draw[arrow] (rng) -- (language);
\draw[arrow] (rng) -- (algorithm);

\end{tikzpicture}
\caption{Diagram klas warstwy Core}
\end{figure}

\newpage

\section{Wzorce Projektowe}

System IO\_RNG wykorzystuje szereg wzorców projektowych zgodnych z zasadami SOLID i Clean Architecture.

\subsection{Clean Architecture (Hexagonal Architecture)}

\textbf{Opis:} Architektura heksagonalna (Porty i Adaptery) oddziela logikę biznesową od szczegółów technicznych.

\textbf{Implementacja w projekcie:}

\begin{tcolorbox}[colback=blue!5!white,colframe=blue!75!black,title=Struktura Clean Architecture]
\begin{itemize}
    \item \textbf{Warstwa Core (Domain):}
    \begin{itemize}
        \item \texttt{entities/} - encje domenowe (RNG, TestResult)
        \item \texttt{interfaces/} - porty (IRNGRepository, IRNGRunner)
        \item \texttt{use\_cases/} - logika biznesowa
    \end{itemize}
    
    \item \textbf{Warstwa Infrastructure (Adapters):}
    \begin{itemize}
        \item \texttt{repositories/} - implementacje repozytoriów
        \item \texttt{runners/} - implementacje runnerów dla różnych języków
        \item \texttt{models.py} - modele Django ORM
    \end{itemize}
    
    \item \textbf{Warstwa API:}
    \begin{itemize}
        \item \texttt{views.py} - ViewSety Django REST Framework
        \item \texttt{serializers.py} - serializatory danych
    \end{itemize}
\end{itemize}
\end{tcolorbox}

\textbf{Zalety:}
\begin{itemize}
    \item Niezależność logiki biznesowej od frameworka
    \item Łatwość testowania (możliwość mockowania interfejsów)
    \item Możliwość łatwej zamiany implementacji (np. zmiana bazy danych)
\end{itemize}

\subsection{Repository Pattern}

\textbf{Cel:} Abstrakcja dostępu do danych, oddzielenie logiki biznesowej od mechanizmów persystencji.

\textbf{Implementacja:}

\begin{lstlisting}[language=Python, caption=Interface repozytorium]
class IRNGRepository(ABC):
    @abstractmethod
    def save(self, rng: RNG) -> RNG:
        pass
    
    @abstractmethod
    def get_by_id(self, rng_id: int) -> Optional[RNG]:
        pass
    
    @abstractmethod
    def get_all(self) -> List[RNG]:
        pass
\end{lstlisting}

\begin{lstlisting}[language=Python, caption=Implementacja dla Django ORM]
class DjangoRNGRepository(IRNGRepository):
    def get_by_id(self, rng_id: int) -> Optional[RNG]:
        try:
            model = RNGModel.objects.get(id=rng_id)
            return self._to_entity(model)
        except RNGModel.DoesNotExist:
            return None
    
    def _to_entity(self, model: RNGModel) -> RNG:
        return RNG(
            id=model.id,
            name=model.name,
            language=Language(model.language),
            algorithm=Algorithm(model.algorithm),
            # ... inne pola
        )
\end{lstlisting}

\textbf{Zalety:}
\begin{itemize}
    \item Enkapsulacja logiki dostępu do danych
    \item Możliwość łatwej zamiany źródła danych (SQLite → PostgreSQL)
    \item Ułatwienie testowania (mock repositories)
\end{itemize}

\newpage

\subsection{Adapter Pattern}

\textbf{Cel:} Umożliwienie współpracy niezgodnych interfejsów.

\textbf{Implementacja - Universal RNG Adapter:}

System używa wzorca Adapter do automatycznego wykrywania i adaptacji różnych implementacji generatorów:

\begin{lstlisting}[language=Python, caption=Universal RNG Adapter]
class UniversalRNGAdapter:
    def __init__(self, module, parameters: Optional[Dict] = None):
        self.module = module
        self.parameters = parameters or {}
        self.generator_func = self._detect_generator_function()
        self.generator_type = self._detect_generator_type()
    
    def _detect_generator_function(self):
        # Priorytet 1: generate()
        if hasattr(self.module, 'generate'):
            return self.module.generate
        
        # Priorytet 2: *_bit_stream()
        for name in dir(self.module):
            if name.endswith('_bit_stream'):
                return getattr(self.module, name)
        
        raise RuntimeError("No generator function found")
    
    def _detect_generator_type(self) -> DataType:
        # Wykrywa typ danych: BITS, INTEGERS, FLOATS
        # ...
\end{lstlisting}

\textbf{Zalety:}
\begin{itemize}
    \item Automatyczne dopasowanie do różnych sygnatur funkcji
    \item Obsługa różnych typów danych wyjściowych (bity, integery, floaty)
    \item Brak konieczności modyfikacji istniejących generatorów
\end{itemize}

\subsection{Strategy Pattern}

\textbf{Cel:} Umożliwienie wyboru algorytmu w czasie wykonania.

\textbf{Implementacja - RNG Runners:}

\begin{lstlisting}[language=Python, caption=Interface strategii]
class IRNGRunner(ABC):
    @abstractmethod
    def can_run(self, rng: RNG) -> bool:
        pass
    
    @abstractmethod
    def generate_numbers(self, rng: RNG, count: int, 
                        seed: int = None) -> List[float]:
        pass
\end{lstlisting}

\begin{lstlisting}[language=Python, caption=Konkretne strategie]
class PythonRNGRunner(IRNGRunner):
    def can_run(self, rng: RNG) -> bool:
        return rng.language == Language.PYTHON
    # ...

class ExeRNGRunner(IRNGRunner):
    def can_run(self, rng: RNG) -> bool:
        return rng.language == Language.EXECUTABLE
    # ...
\end{lstlisting}

\textbf{Wybór strategii w runtime:}
\begin{lstlisting}[language=Python]
def _find_runner(self, rng: RNG) -> IRNGRunner:
    for runner in self.runners:
        if runner.can_run(rng):
            return runner
    raise RuntimeError(f"No runner for {rng.language}")
\end{lstlisting}

\newpage

\subsection{Use Case Pattern}

\textbf{Cel:} Enkapsulacja logiki biznesowej w pojedynczych, specyficznych przypadkach użycia.

\textbf{Implementacja:}

\begin{lstlisting}[language=Python, caption=Run RNG Test Use Case]
class RunRNGTestUseCase:
    def __init__(self,
                 rng_repository: IRNGRepository,
                 result_repository: ITestResultRepository,
                 runners: List[IRNGRunner]):
        self.rng_repository = rng_repository
        self.result_repository = result_repository
        self.runners = runners
    
    def execute(self, rng_id: int, test_name: str, 
                samples_count: int, seed: int = None) -> TestResult:
        # 1. Pobierz RNG z repozytorium
        rng = self.rng_repository.get_by_id(rng_id)
        
        # 2. Znajdz odpowiedni runner
        runner = self._find_runner(rng)
        
        # 3. Generuj liczby losowe
        numbers = runner.generate_numbers(rng, samples_count, seed)
        
        # 4. Wykonaj test statystyczny
        test_result = self._perform_statistical_test(...)
        
        # 5. Zapisz wynik
        return self.result_repository.save(result)
\end{lstlisting}

\textbf{Inne Use Cases w systemie:}
\begin{itemize}
    \item \texttt{CompareRNGsUseCase} - porównywanie generatorów
    \item \texttt{GetTestResultsUseCase} - pobieranie wyników testów
    \item \texttt{TestCustomBitsUseCase} - testowanie własnych sekwencji bitów
\end{itemize}

\subsection{Dependency Injection}

\textbf{Cel:} Odwrócenie kontroli, zwiększenie testowalności.

\textbf{Implementacja w ViewSet:}

\begin{lstlisting}[language=Python, caption=Wstrzykiwanie zależności]
class RNGViewSet(viewsets.ViewSet):
    def __init__(self, **kwargs):
        super().__init__(**kwargs)
        # Tworzenie konkretnych implementacji
        self.rng_repository = DjangoRNGRepository()
        self.result_repository = DjangoTestResultRepository()
        self.runners = [
            PythonRNGRunner(),
            ExeRNGRunner()
        ]
    
    @action(detail=True, methods=['post'])
    def run_test(self, request, pk=None):
        # Wstrzykiwanie do Use Case
        use_case = RunRNGTestUseCase(
            self.rng_repository,
            self.result_repository,
            self.runners
        )
        result = use_case.execute(...)
        return Response(...)
\end{lstlisting}

\newpage

\subsection{Factory Pattern (Implicit)}

System używa niejawnego wzorca Factory w procesie tworzenia encji z modeli Django:

\begin{lstlisting}[language=Python, caption=Konwersja Model do Entity]
def _to_entity(self, model: RNGModel) -> RNG:
    """Factory method - creates entity from Django model"""
    return RNG(
        id=model.id,
        name=model.name,
        language=Language(model.language),
        algorithm=Algorithm(model.algorithm),
        description=model.description,
        code_path=model.code_path,
        parameters=model.parameters or {},
        is_active=model.is_active
    )
\end{lstlisting}

\section{Zasady Programowania Obiektowego}

\subsection{SOLID Principles}

System IO\_RNG konsekwentnie stosuje zasady SOLID:

\subsubsection{Single Responsibility Principle (SRP)}

Każda klasa ma jedną, dobrze zdefiniowaną odpowiedzialność:

\begin{tcolorbox}[colback=green!5!white,colframe=green!75!black]
\textbf{Przykłady:}
\begin{itemize}
    \item \texttt{RNG} - reprezentuje generator (encja domeny)
    \item \texttt{DjangoRNGRepository} - zarządza persystencją RNG
    \item \texttt{PythonRNGRunner} - wykonuje generatory Python
    \item \texttt{RunRNGTestUseCase} - orkiestruje proces testowania
    \item \texttt{UniversalRNGAdapter} - adaptuje różne formaty generatorów
\end{itemize}
\end{tcolorbox}

\subsubsection{Open/Closed Principle (OCP)}

Klasy są otwarte na rozszerzenia, zamknięte na modyfikacje:

\begin{lstlisting}[language=Python, caption=Przyklad OCP - dodanie nowego runnera]
# Adding support for new language without modifying existing code
class RustRNGRunner(IRNGRunner):
    def can_run(self, rng: RNG) -> bool:
        return rng.language == Language.RUST
    
    def generate_numbers(self, rng: RNG, ...) -> List[float]:
        # Rust implementation
        pass

# In ViewSet - just add to list
self.runners = [
    PythonRNGRunner(),
    ExeRNGRunner(),
    RustRNGRunner()  # New runner
]
\end{lstlisting}

\subsubsection{Liskov Substitution Principle (LSP)}

Wszystkie implementacje interfejsów są podstawialne:

\begin{lstlisting}[language=Python]
# Every IRNGRunner implementation is substitutable
def test_with_any_runner(runner: IRNGRunner, rng: RNG):
    if runner.can_run(rng):
        numbers = runner.generate_numbers(rng, 1000)
        # Works with PythonRNGRunner, ExeRNGRunner, etc.
\end{lstlisting}

\subsubsection{Interface Segregation Principle (ISP)}

Interfejsy są minimalistyczne i skoncentrowane:

\begin{itemize}
    \item \texttt{IRNGRepository} - tylko operacje na RNG
    \item \texttt{ITestResultRepository} - tylko operacje na wynikach
    \item \texttt{IRNGRunner} - tylko generowanie liczb
\end{itemize}

\subsubsection{Dependency Inversion Principle (DIP)}

Moduły wysokopoziomowe zależą od abstrakcji, nie od konkretnych implementacji:

\begin{lstlisting}[language=Python]
class RunRNGTestUseCase:
    # Depends on abstractions (interfaces)
    def __init__(self,
                 rng_repository: IRNGRepository,  # Abstraction
                 result_repository: ITestResultRepository,  # Abstraction
                 runners: List[IRNGRunner]):  # Abstraction
        # ...
\end{lstlisting}

\newpage

\subsection{Enkapsulacja}

\textbf{Enkapsulacja danych i logiki w encjach:}

\begin{lstlisting}[language=Python, caption=Enkapsulacja w RNG Entity]
@dataclass
class RNG:
    # Public attributes (dataclass)
    name: str
    language: Language
    algorithm: Algorithm
    # ...
    
    def __post_init__(self):
        """Private validation - called automatically"""
        if not self.name:
            raise ValueError("RNG name cannot be empty")
        if not self.code_path:
            raise ValueError("Code path cannot be empty")
    
    def validate_for_execution(self) -> bool:
        """Public business method"""
        return self.is_active and bool(self.code_path)
    
    def get_parameter(self, key: str, default: Any = None) -> Any:
        """Safe access to parameters"""
        return self.parameters.get(key, default)
\end{lstlisting}

\subsection{Polimorfizm}

\textbf{Polimorfizm w runnerach:}

\begin{lstlisting}[language=Python, caption=Rozne implementacje tego samego interfejsu]
# Common interface
numbers = runner.generate_numbers(rng, count, seed)

# PythonRNGRunner - loads Python module and executes
# ExeRNGRunner - runs .exe file and parses stdout
# CppRNGRunner - compiles and runs C++ code
# All return List[float]
\end{lstlisting}

\subsection{Abstrakcja}

System używa abstrakcji na wielu poziomach:

\begin{itemize}
    \item \textbf{Encje} - abstrakcja koncepcji biznesowych (RNG, TestResult)
    \item \textbf{Interfejsy} - abstrakcja operacji (IRNGRunner, IRNGRepository)
    \item \textbf{Enumy} - abstrakcja wartości (Language, Algorithm, DataType)
\end{itemize}

\section{Struktury Danych}

\subsection{Encje Domenowe (Dataclasses)}

System używa Python dataclasses jako immutable value objects:

\begin{lstlisting}[language=Python, caption=Struktura encji RNG]
@dataclass
class RNG:
    name: str
    language: Language
    algorithm: Algorithm
    description: str
    code_path: str
    parameters: Optional[Dict[str, Any]] = None
    id: Optional[int] = None
    is_active: bool = True
\end{lstlisting}

\textbf{Zalety dataclasses:}
\begin{itemize}
    \item Automatyczna generacja \texttt{\_\_init\_\_}, \texttt{\_\_repr\_\_}, \texttt{\_\_eq\_\_}
    \item Type hints dla wszystkich pól
    \item Niezmienność (immutability) z \texttt{frozen=True}
    \item Walidacja w \texttt{\_\_post\_init\_\_}
\end{itemize}

\subsection{Enumy (Typy wyliczeniowe)}

System definiuje silnie typowane enumy dla wartości biznesowych:

\begin{lstlisting}[language=Python, caption=Enumy domenowe]
class Language(Enum):
    PYTHON = "python"
    JAVA = "java"
    CPP = "cpp"
    RUST = "rust"
    EXECUTABLE = "executable"

class Algorithm(Enum):
    LINEAR_CONGRUENTIAL = "linear_congruential"
    MERSENNE_TWISTER = "mersenne_twister"
    XORSHIFT = "xorshift"
    PCG = "pcg"
    SYSTEM = "system"

class DataType(Enum):
    BITS = "bits"
    INTEGERS = "integers"
    FLOATS = "floats"
\end{lstlisting}

\textbf{Zalety enumów:}
\begin{itemize}
    \item Bezpieczeństwo typów
    \item Niemożliwość użycia nieprawidłowych wartości
    \item Autocomplete w IDE
    \item Czytelność kodu
\end{itemize}

\subsection{Słowniki (Dict) jako Parametry}

System używa słowników JSON do przechowywania dynamicznych parametrów:

\begin{lstlisting}[language=Python, caption=Parametry generatorow]
# Example LCG parameters
parameters = {
    "a": 1103515245,      # multiplier
    "c": 12345,           # increment
    "m": 2**31,           # modulus
    "bits_per_value": 31  # number of bits
}

# Safe access via method
def get_parameter(self, key: str, default: Any = None) -> Any:
    return self.parameters.get(key, default) if self.parameters else default
\end{lstlisting}

\subsection{Listy typowane}

System używa list z type hints dla bezpieczeństwa typów:

\begin{lstlisting}[language=Python]
# List of bits
bits: List[int] = [0, 1, 0, 1, 1, 0, ...]

# List of floats
numbers: List[float] = [0.234, 0.891, 0.456, ...]

# List of entities
rngs: List[RNG] = repository.get_all()

# List of runners
runners: List[IRNGRunner] = [PythonRNGRunner(), ExeRNGRunner()]
\end{lstlisting}

\subsection{Tuple dla zwracania wielu wartości}

\begin{lstlisting}[language=Python, caption=Tuple z type hints]
def generate_raw(self, rng: RNG, count: int, 
                 seed: int = None) -> Tuple[List[Any], DataType]:
    """Returns (data, data_type)"""
    # ...
    return (data, DataType.BITS)
\end{lstlisting}

\subsection{Optional dla wartości nullable}

\begin{lstlisting}[language=Python]
def get_by_id(self, rng_id: int) -> Optional[RNG]:
    """Can return RNG or None"""
    try:
        model = RNGModel.objects.get(id=rng_id)
        return self._to_entity(model)
    except RNGModel.DoesNotExist:
        return None
\end{lstlisting}

\subsection{Struktury danych w bazie}

\subsubsection{Model RNG (Tabela rngs)}

\begin{lstlisting}[language=Python, caption=Django Model]
class RNGModel(models.Model):
    name = models.CharField(max_length=200)
    language = models.CharField(max_length=50)
    algorithm = models.CharField(max_length=100)
    description = models.TextField()
    code_path = models.CharField(max_length=500)
    parameters = models.JSONField(null=True, blank=True)
    is_active = models.BooleanField(default=True)
    created_at = models.DateTimeField(auto_now_add=True)
    
    class Meta:
        db_table = 'rngs'
        ordering = ['-created_at']
\end{lstlisting}

\textbf{Indeksy:}
\begin{itemize}
    \item Primary key na \texttt{id}
    \item Ordering index na \texttt{created\_at}
\end{itemize}

\subsubsection{Model TestResult (Tabela test\_results)}

\begin{lstlisting}[language=Python]
class TestResultModel(models.Model):
    rng = models.ForeignKey(RNGModel, on_delete=models.CASCADE,
                           related_name='test_results')
    test_name = models.CharField(max_length=100)
    passed = models.BooleanField()
    score = models.FloatField()
    execution_time_ms = models.FloatField()
    samples_count = models.IntegerField()
    statistics = models.JSONField()
    error_message = models.TextField(null=True, blank=True)
    created_at = models.DateTimeField(auto_now_add=True)
    
    class Meta:
        db_table = 'test_results'
        ordering = ['-created_at']
        indexes = [
            models.Index(fields=['rng', '-created_at']),
            models.Index(fields=['test_name', '-created_at']),
        ]
\end{lstlisting}

\textbf{Indeksy optymalizacyjne:}
\begin{itemize}
    \item Composite index: \texttt{(rng\_id, created\_at)} - szybkie pobieranie wyników dla RNG
    \item Composite index: \texttt{(test\_name, created\_at)} - filtrowanie po typie testu
\end{itemize}

\newpage

\section{Diagram Przepływu Danych}

\subsection{Proces uruchomienia testu RNG}

\begin{figure}[H]
\centering
\begin{tikzpicture}[
    node distance=1.2cm,
    process/.style={rectangle, rounded corners, draw=black!80, fill=blue!20, thick, minimum width=3cm, minimum height=1cm, align=center},
    decision/.style={diamond, draw=black!80, fill=yellow!20, thick, minimum width=2.5cm, minimum height=1cm, align=center, aspect=2},
    data/.style={trapezium, trapezium left angle=70, trapezium right angle=110, draw=black!80, fill=green!20, thick, minimum width=2.5cm, minimum height=0.8cm, align=center},
    arrow/.style={thick,->,>=stealth}
]

\node[data] (request) {POST /api/rngs/\{id\}/run\_test};
\node[process, below=of request] (viewset) {RNGViewSet\\run\_test()};
\node[process, below=of viewset] (usecase) {RunRNGTestUseCase\\execute()};
\node[process, below=of usecase] (getrnq) {Repository\\get\_by\_id()};
\node[decision, below=of getrnq] (exists) {RNG\\exists?};
\node[process, below=of exists, yshift=-0.5cm] (findrunner) {Find appropriate\\runner};
\node[decision, below=of findrunner] (canrun) {Runner\\found?};
\node[process, below=of canrun, yshift=-0.5cm] (generate) {Runner\\generate\_numbers()};
\node[process, below=of generate] (test) {Perform\\statistical test};
\node[process, below=of test] (save) {Repository\\save(result)};
\node[data, below=of save] (response) {HTTP 200\\TestResult JSON};

\node[data, right=3cm of exists] (error1) {HTTP 404\\Not Found};
\node[data, right=3cm of canrun] (error2) {HTTP 500\\No Runner};

\draw[arrow] (request) -- (viewset);
\draw[arrow] (viewset) -- (usecase);
\draw[arrow] (usecase) -- (getrnq);
\draw[arrow] (getrnq) -- (exists);
\draw[arrow] (exists) -- node[left] {yes} (findrunner);
\draw[arrow] (exists) -- node[above] {no} (error1);
\draw[arrow] (findrunner) -- (canrun);
\draw[arrow] (canrun) -- node[left] {yes} (generate);
\draw[arrow] (canrun) -- node[above] {no} (error2);
\draw[arrow] (generate) -- (test);
\draw[arrow] (test) -- (save);
\draw[arrow] (save) -- (response);

\end{tikzpicture}
\caption{Sekwencja uruchomienia testu RNG}
\end{figure}

\newpage

\subsection{Proces generowania liczb - Python Runner}

\begin{figure}[H]
\centering
\begin{tikzpicture}[
    node distance=1.2cm,
    process/.style={rectangle, rounded corners, draw=black!80, fill=blue!20, thick, minimum width=3.5cm, minimum height=0.9cm, align=center},
    decision/.style={diamond, draw=black!80, fill=yellow!20, thick, minimum width=2.2cm, minimum height=0.9cm, align=center, aspect=2},
    arrow/.style={thick,->,>=stealth}
]

\node[process] (start) {PythonRNGRunner\\generate\_numbers()};
\node[process, below=of start] (load) {Load Python module\\importlib};
\node[process, below=of load] (adapter) {UniversalRNGAdapter\\create};
\node[process, below=of adapter] (detect) {Detect generator\\function};
\node[decision, below=of detect] (functype) {Function\\type?};

\node[process, below left=1cm and 1cm of functype] (bits) {*\_bit\_stream()\\generate bits};
\node[process, below right=1cm and 1cm of functype] (floats) {generate()\\generate floats};

\node[process, below=2cm of functype] (convert) {Convert to\\List[float]};
\node[process, below=of convert] (return) {Return numbers};

\draw[arrow] (start) -- (load);
\draw[arrow] (load) -- (adapter);
\draw[arrow] (adapter) -- (detect);
\draw[arrow] (detect) -- (functype);
\draw[arrow] (functype) -- node[left, xshift=-0.2cm] {bits} (bits);
\draw[arrow] (functype) -- node[right, xshift=0.2cm] {floats} (floats);
\draw[arrow] (bits) -- (convert);
\draw[arrow] (floats) -- (convert);
\draw[arrow] (convert) -- (return);

\end{tikzpicture}
\caption{Proces generowania liczb w Python Runner}
\end{figure}

\section{Szczegóły Implementacji}

\subsection{Kompresja danych - Base64}

System używa kompresji bitów do Base64 dla oszczędności transferu danych:

\begin{lstlisting}[language=Python, caption=Kompresja bitow]
def compress_bits_to_base64(bits: List[int]) -> str:
    """
    Packs bits list [0,1,0,1,...] to base64.
    Savings: ~8x (3 bytes/bit to 0.375 bytes/bit)
    """
    byte_array = bytearray()
    
    for i in range(0, len(bits), 8):
        byte_val = 0
        for j in range(8):
            if i + j < len(bits):
                # MSB first
                byte_val |= (bits[i + j] << (7 - j))
        byte_array.append(byte_val)
    
    return base64.b64encode(byte_array).decode('ascii')
\end{lstlisting}

\textbf{Przykład:}
\begin{itemize}
    \item Bez kompresji: \texttt{[1,0,1,0,1,0,1,0]} → 24 bajty JSON
    \item Z kompresją: \texttt{"qg=="} → 4 bajty
    \item Oszczędność: 83\%
\end{itemize}

\subsection{Testy Statystyczne}

System implementuje szereg testów statystycznych:

\begin{enumerate}
    \item \textbf{Chi-Square Test} - test zgodności rozkładu
    \item \textbf{Kolmogorov-Smirnov Test} - test równomierności
    \item \textbf{Runs Test} - test niezależności sekwencji
    \item \textbf{Spectral Test (FFT)} - test częstotliwościowy
    \item \textbf{Poker Test} - test sekwencji 5-bitowych
    \item \textbf{Monobit Test} - test proporcji bitów
    \item \textbf{Serial Test} - test par bitów
    \item \textbf{Autocorrelation Test} - test autokorelacji
\end{enumerate}

\subsection{Obsługa różnych języków programowania}

\subsubsection{Python Runner}

\begin{itemize}
    \item Dynamiczne ładowanie modułów: \texttt{importlib.util}
    \item Automatyczna detekcja funkcji generatora
    \item Wsparcie dla parametryzowanych generatorów (LCG, AWCG)
\end{itemize}

\subsubsection{Executable Runner}

\begin{itemize}
    \item Uruchamianie prekompilowanych plików .exe
    \item Komunikacja przez stdout/stdin
    \item Parsowanie wyjścia tekstowego
    \item Wsparcie dla Rust, C++, C\#
\end{itemize}

\newpage

\section{Architektura Backendowa}

\subsection{Przegląd struktury projektu Django}

Backend systemu IO\_RNG zbudowany jest w oparciu o framework Django REST Framework, stosując zasady Clean Architecture. Projekt podzielony jest na wyraźnie oddzielone warstwy, co zapewnia separację odpowiedzialności i łatwość testowania.

\begin{tcolorbox}[colback=blue!5!white,colframe=blue!75!black,title=Struktura katalogów backendu]
\small
\begin{verbatim}
backend/
  io_rng/                      # Glowna aplikacja Django
    core/                      # Warstwa domeny (Core Layer)
      entities/                # Encje domenowe
        rng.py                 # RNG (Generator)
        test_result.py         # Wynik testu
        enums.py               # Language, Algorithm, DataType
      interfaces/              # Interfejsy (ABC)
        repository.py          # IRNGRepository, ITestResultRepository
        runner.py              # IRNGRunner
      use_cases/               # Logika biznesowa
        run_rng_test.py        # Uruchomienie testu
        compare_rngs.py        # Porownanie generatorow
        get_test_results.py    # Pobieranie wynikow
        test_custom_bits.py    # Test wlasnych bitow
    
    infrastructure/            # Warstwa infrastruktury
      models.py                # Django ORM Models
      repositories/            # Implementacje repozytoriow
        django_repositories.py
      runners/                 # Implementacje runnerow
        python_runner.py       # Python RNG Runner
        exe_runner.py          # Executable Runner
        universal_adapter.py   # Uniwersalny adapter
    
    api/                       # Warstwa API
      views.py                 # Django REST Framework ViewSets
      serializers.py           # DRF Serializers
      urls.py                  # Routing API
    
    settings.py                # Konfiguracja Django
    urls.py                    # Glowny routing
    wsgi.py                    # WSGI application
  
  db.sqlite3                   # Baza danych SQLite
  manage.py                    # Django CLI
  requirements.txt             # Zaleznosci Python
\end{verbatim}
\end{tcolorbox}

\subsection{Core Layer - Warstwa Domeny}

Warstwa Core zawiera całą logikę biznesową systemu i jest całkowicie niezależna od frameworka Django. Wszystkie komponenty w tej warstwie operują na czystych obiektach Pythona (dataclasses, enums, abstrakcyjne klasy).

\subsubsection{Encje Domenowe}

\textbf{RNG Entity} - reprezentuje generator liczb pseudolosowych:

\begin{lstlisting}[language=Python, caption=Encja RNG]
@dataclass
class RNG:
    """Generator liczb pseudolosowych - encja domenowa"""
    name: str                           # Nazwa generatora
    language: Language                  # Język implementacji
    algorithm: Algorithm                # Algorytm RNG
    description: str                    # Opis generatora
    code_path: str                      # Ścieżka do pliku
    parameters: Optional[Dict] = None   # Parametry (LCG: a,c,m)
    id: Optional[int] = None            # ID w bazie
    is_active: bool = True              # Czy aktywny

    def __post_init__(self):
        """Walidacja danych po inicjalizacji"""
        if not self.name:
            raise ValueError("RNG name cannot be empty")
        if not self.code_path:
            raise ValueError("Code path cannot be empty")

    def validate_for_execution(self) -> bool:
        """Sprawdza czy generator jest gotowy do wykonania"""
        return self.is_active and bool(self.code_path)

    def get_parameter(self, key: str, default: Any = None) -> Any:
        """Bezpieczny dostęp do parametrów"""
        return self.parameters.get(key, default) if self.parameters else default
\end{lstlisting}

\textbf{TestResult Entity} - reprezentuje wynik testu statystycznego:

\begin{lstlisting}[language=Python, caption=Encja TestResult]
@dataclass
class TestResult:
    """Wynik testu statystycznego - encja domenowa"""
    rng_id: int                         # ID testowanego RNG
    test_name: str                      # Nazwa testu
    passed: bool                        # Czy test przeszedł
    score: float                        # Wynik (0-1)
    execution_time_ms: float            # Czas wykonania
    samples_count: int                  # Liczba próbek
    statistics: Dict[str, Any]          # Szczegółowe statystyki
    seed: Optional[int] = None          # Seed użyty do testu
    test_parameters: Optional[Dict] = None  # Parametry testu
    error_message: Optional[str] = None # Komunikat błędu
    id: Optional[int] = None            # ID w bazie

    def __post_init__(self):
        """Walidacja wyniku testu"""
        if not 0.0 <= self.score <= 1.0:
            raise ValueError("Score must be between 0 and 1")
        if self.execution_time_ms < 0:
            raise ValueError("Execution time cannot be negative")
\end{lstlisting}

\subsubsection{Interfejsy (Porty)}

System definiuje abstrakcyjne interfejsy, które muszą być implementowane przez warstwę Infrastructure:

\begin{lstlisting}[language=Python, caption=Interfejsy repozytoriów]
class IRNGRepository(ABC):
    """Interface repozytorium generatorów RNG"""

    @abstractmethod
    def save(self, rng: RNG) -> RNG:
        """Zapisuje generator (create lub update)"""
        pass

    @abstractmethod
    def get_by_id(self, rng_id: int) -> Optional[RNG]:
        """Pobiera generator po ID"""
        pass

    @abstractmethod
    def get_all(self) -> List[RNG]:
        """Pobiera wszystkie generatory"""
        pass

    @abstractmethod
    def delete(self, rng_id: int) -> bool:
        """Usuwa generator"""
        pass

class ITestResultRepository(ABC):
    """Interface repozytorium wyników testów"""

    @abstractmethod
    def save(self, result: TestResult) -> TestResult:
        pass

    @abstractmethod
    def get_by_rng(self, rng_id: int) -> List[TestResult]:
        pass

    @abstractmethod
    def get_all(self) -> List[TestResult]:
        pass
\end{lstlisting}

\begin{lstlisting}[language=Python, caption=Interface runnera]
class IRNGRunner(ABC):
    """Interface runnera dla różnych języków programowania"""

    @abstractmethod
    def can_run(self, rng: RNG) -> bool:
        """Sprawdza czy runner może uruchomić dany RNG"""
        pass

    @abstractmethod
    def generate_numbers(self, rng: RNG, count: int,
                        seed: int = None) -> List[float]:
        """Generuje liczby losowe z zakresu [0,1]"""
        pass

    @abstractmethod
    def generate_raw(self, rng: RNG, count: int,
                    seed: int = None) -> Tuple[List[Any], DataType]:
        """Generuje surowe dane (bity, integery lub floaty)"""
        pass
\end{lstlisting}

\subsubsection{Use Cases - Logika Biznesowa}

Use cases orkiestrują przepływ danych między warstwami i implementują logikę biznesową:

\begin{lstlisting}[language=Python, caption=RunRNGTestUseCase - główny przypadek użycia]
class RunRNGTestUseCase:
    """Przypadek użycia: Uruchomienie testu statystycznego na RNG"""

    def __init__(self,
                 rng_repository: IRNGRepository,
                 result_repository: ITestResultRepository,
                 runners: List[IRNGRunner]):
        # Dependency Injection - zależności wstrzykiwane przez konstruktor
        self.rng_repository = rng_repository
        self.result_repository = result_repository
        self.runners = runners

    def execute(self, rng_id: int, test_name: str,
                samples_count: int, seed: int = None,
                test_parameters: Dict = None) -> TestResult:
        """
        Wykonuje test statystyczny:
        1. Pobiera RNG z repozytorium
        2. Znajduje odpowiedni runner
        3. Generuje liczby/bity
        4. Przeprowadza test statystyczny
        5. Zapisuje wynik
        """

        # 1. Pobranie generatora
        rng = self.rng_repository.get_by_id(rng_id)
        if not rng:
            raise ValueError(f"RNG with id {rng_id} not found")

        # 2. Znalezienie odpowiedniego runnera
        runner = self._find_runner(rng)

        # 3. Generowanie danych
        start_time = time.time()

        # Testy NIST wymagają bitów, inne testy floatów
        if test_name.startswith('nist_') or test_name.startswith('diehard_'):
            bits, data_type = runner.generate_raw(rng, samples_count, seed)
            # Konwersja do bitów jeśli potrzeba
            if data_type != DataType.BITS:
                bits = self._convert_to_bits(bits, data_type)
        else:
            # Podstawowe testy (frequency, uniformity) używają floatów
            numbers = runner.generate_numbers(rng, samples_count, seed)

        generation_time = (time.time() - start_time) * 1000

        # 4. Przeprowadzenie testu statystycznego
        test_result = self._perform_test(
            test_name,
            bits if 'bits' in locals() else numbers,
            test_parameters
        )

        # 5. Utworzenie i zapisanie wyniku
        result = TestResult(
            rng_id=rng_id,
            test_name=test_name,
            passed=test_result['passed'],
            score=test_result['score'],
            execution_time_ms=generation_time + test_result['test_time_ms'],
            samples_count=samples_count,
            statistics=test_result['statistics'],
            seed=seed,
            test_parameters=test_parameters
        )

        return self.result_repository.save(result)
\end{lstlisting}

\newpage

\subsection{Infrastructure Layer - Warstwa Infrastruktury}

Warstwa Infrastructure zawiera konkretne implementacje interfejsów zdefiniowanych w Core Layer. Używa Django ORM do persystencji danych.

\subsubsection{Django Models}

\begin{lstlisting}[language=Python, caption=Modele Django ORM]
class RNGModel(models.Model):
    """Model Django dla generatorów RNG"""

    name = models.CharField(max_length=200, unique=True)
    language = models.CharField(max_length=50)  # python, cpp, rust, etc.
    algorithm = models.CharField(max_length=100)
    description = models.TextField(blank=True)
    code_path = models.CharField(max_length=500)  # Względna ścieżka
    parameters = models.JSONField(null=True, blank=True)  # Parametry JSON
    is_active = models.BooleanField(default=True)
    created_at = models.DateTimeField(auto_now_add=True)
    updated_at = models.DateTimeField(auto_now=True)

    class Meta:
        db_table = 'rngs'
        ordering = ['-created_at']
        indexes = [
            models.Index(fields=['language']),
            models.Index(fields=['algorithm']),
        ]

class TestResultModel(models.Model):
    """Model Django dla wyników testów"""

    rng = models.ForeignKey(RNGModel, on_delete=models.CASCADE,
                           related_name='test_results')
    test_name = models.CharField(max_length=100)
    passed = models.BooleanField()
    score = models.FloatField()  # 0.0 - 1.0
    execution_time_ms = models.FloatField()
    samples_count = models.IntegerField()
    statistics = models.JSONField()  # Szczegółowe statystyki
    seed = models.IntegerField(null=True, blank=True)
    test_parameters = models.JSONField(null=True, blank=True)
    error_message = models.TextField(null=True, blank=True)
    created_at = models.DateTimeField(auto_now_add=True)

    class Meta:
        db_table = 'test_results'
        ordering = ['-created_at']
        indexes = [
            models.Index(fields=['rng', '-created_at']),
            models.Index(fields=['test_name', '-created_at']),
            models.Index(fields=['passed']),
        ]
\end{lstlisting}

\subsubsection{Repository Implementations}

Repozytoria implementują wzorzec Repository, przekształcając encje domenowe w modele Django i vice versa:

\begin{lstlisting}[language=Python, caption=Implementacja DjangoRNGRepository]
class DjangoRNGRepository(IRNGRepository):
    """Konkretna implementacja repozytorium używająca Django ORM"""

    def save(self, rng: RNG) -> RNG:
        """Zapisuje lub aktualizuje RNG"""
        if rng.id:
            # Update istniejącego
            model = RNGModel.objects.get(id=rng.id)
            self._update_model_from_entity(model, rng)
        else:
            # Utworzenie nowego
            model = self._create_model_from_entity(rng)

        model.save()
        return self._to_entity(model)

    def get_by_id(self, rng_id: int) -> Optional[RNG]:
        """Pobiera RNG po ID, zwraca None jeśli nie istnieje"""
        try:
            model = RNGModel.objects.get(id=rng_id)
            return self._to_entity(model)
        except RNGModel.DoesNotExist:
            return None

    def _to_entity(self, model: RNGModel) -> RNG:
        """Konwersja Django Model -> Domain Entity (Factory Method)"""
        return RNG(
            id=model.id,
            name=model.name,
            language=Language(model.language),
            algorithm=Algorithm(model.algorithm),
            description=model.description,
            code_path=model.code_path,
            parameters=model.parameters or {},
            is_active=model.is_active
        )

    def _create_model_from_entity(self, rng: RNG) -> RNGModel:
        """Konwersja Domain Entity -> Django Model"""
        return RNGModel(
            name=rng.name,
            language=rng.language.value,
            algorithm=rng.algorithm.value,
            description=rng.description,
            code_path=rng.code_path,
            parameters=rng.parameters,
            is_active=rng.is_active
        )
\end{lstlisting}

\newpage

\subsubsection{System Runnerów}

System runnerów umożliwia wykonywanie generatorów napisanych w różnych językach programowania. Każdy runner implementuje interfejs \texttt{IRNGRunner}.

\begin{tcolorbox}[colback=green!5!white,colframe=green!75!black,title=Obsługiwane języki]
\begin{itemize}
    \item \textbf{Python} - dynamiczne ładowanie modułów (importlib)
    \item \textbf{C++} - kompilacja i uruchomienie plików .exe
    \item \textbf{Rust} - uruchomienie prekompilowanych binarek
    \item \textbf{C\#} - uruchomienie plików .exe/.dll
    \item \textbf{Java} - kompilacja i uruchomienie .class (TODO)
\end{itemize}
\end{tcolorbox}

\textbf{Python RNG Runner:}

\begin{lstlisting}[language=Python, caption=PythonRNGRunner - wykonywanie generatorów Python]
class PythonRNGRunner(IRNGRunner):
    """Runner dla generatorów napisanych w Pythonie"""

    BASE_PATH = Path(__file__).parent.parent.parent / "algorytmy"

    def can_run(self, rng: RNG) -> bool:
        """Sprawdza czy runner obsługuje dany RNG"""
        return rng.language == Language.PYTHON

    def generate_numbers(self, rng: RNG, count: int,
                        seed: int = None) -> List[float]:
        """Generuje liczby float [0,1]"""
        # 1. Załaduj moduł Python
        module = self._load_module(rng.code_path)

        # 2. Utwórz Universal Adapter
        adapter = UniversalRNGAdapter(module, rng.parameters)

        # 3. Wygeneruj surowe dane
        raw_data, data_type = adapter.generate_raw(count, seed)

        # 4. Konwertuj do floatów [0,1]
        return self._convert_to_floats(raw_data, data_type)

    def _load_module(self, code_path: str):
        """Dynamicznie ładuje moduł Python używając importlib"""
        full_path = self.BASE_PATH / code_path

        spec = importlib.util.spec_from_file_location(
            "rng_module", full_path
        )
        module = importlib.util.module_from_spec(spec)
        spec.loader.exec_module(module)

        return module
\end{lstlisting}

\textbf{Executable RNG Runner:}

\begin{lstlisting}[language=Python, caption=ExeRNGRunner - uruchamianie plików wykonywalnych]
class ExeRNGRunner(IRNGRunner):
    """Runner dla prekompilowanych plików wykonywalnych (C++, Rust, C#)"""

    def can_run(self, rng: RNG) -> bool:
        return rng.language == Language.EXECUTABLE

    def generate_raw(self, rng: RNG, count: int,
                    seed: int = None) -> Tuple[List[Any], DataType]:
        """Uruchamia plik .exe i parsuje stdout"""

        # 1. Znajdź plik wykonywalny
        exe_path = self._resolve_executable_path(rng.code_path)

        # 2. Przygotuj argumenty
        args = [str(exe_path), str(count)]
        if seed is not None:
            args.append(str(seed))

        # 3. Uruchom proces
        result = subprocess.run(
            args,
            capture_output=True,
            text=True,
            timeout=30  # 30s timeout
        )

        if result.returncode != 0:
            raise RuntimeError(f"Executable failed: {result.stderr}")

        # 4. Parsuj wyjście (format: liczba na linię lub JSON)
        output = result.stdout.strip()

        # Wykryj format wyjścia
        if output.startswith('[') or output.startswith('{'):
            # Format JSON
            data = json.loads(output)
            return data['bits'], DataType.BITS
        else:
            # Format tekstowy - liczby w liniach
            numbers = [float(line) for line in output.split('\n') if line]
            return numbers, DataType.FLOATS
\end{lstlisting}

\textbf{Universal RNG Adapter:}

Adapter automatycznie wykrywa typ funkcji generatora i dostosowuje się do jej sygnatury:

\begin{lstlisting}[language=Python, caption=UniversalRNGAdapter - automatyczna detekcja]
class UniversalRNGAdapter:
    """
    Uniwersalny adapter dla różnych typów generatorów.
    Automatycznie wykrywa:
    - Typ funkcji (generate, *_bit_stream, etc.)
    - Typ danych wyjściowych (bits, integers, floats)
    - Wymagane parametry
    """

    def __init__(self, module, parameters: Optional[Dict] = None):
        self.module = module
        self.parameters = parameters or {}
        self.generator_func = self._detect_generator_function()
        self.data_type = self._detect_data_type()

    def _detect_generator_function(self):
        """Wykrywa funkcję generatora w module"""
        # Priorytet 1: funkcja generate()
        if hasattr(self.module, 'generate'):
            return self.module.generate

        # Priorytet 2: funkcje *_bit_stream()
        for name in dir(self.module):
            if name.endswith('_bit_stream'):
                return getattr(self.module, name)

        # Priorytet 3: pierwsza funkcja publiczna
        for name in dir(self.module):
            if not name.startswith('_'):
                attr = getattr(self.module, name)
                if callable(attr):
                    return attr

        raise RuntimeError("No generator function found in module")

    def _detect_data_type(self) -> DataType:
        """Wykrywa typ danych zwracanych przez generator"""
        # Sprawdź nazwę funkcji
        func_name = self.generator_func.__name__

        if 'bit' in func_name.lower():
            return DataType.BITS
        elif 'int' in func_name.lower():
            return DataType.INTEGERS
        else:
            return DataType.FLOATS

    def generate_raw(self, count: int, seed: int = None) -> Tuple[List, DataType]:
        """Generuje surowe dane odpowiadając na sygnaturę funkcji"""

        # Sprawdź sygnaturę funkcji
        sig = inspect.signature(self.generator_func)

        # Przygotuj argumenty
        kwargs = {}

        if 'count' in sig.parameters or 'n' in sig.parameters:
            kwargs['count'] = count
        if 'seed' in sig.parameters:
            kwargs['seed'] = seed

        # Dodaj parametry generatora (dla LCG: a, c, m)
        for param_name in sig.parameters:
            if param_name in self.parameters:
                kwargs[param_name] = self.parameters[param_name]

        # Wywołaj funkcję
        data = self.generator_func(**kwargs)

        return list(data), self.data_type
\end{lstlisting}

\newpage

\subsection{API Layer - Warstwa Prezentacji}

Warstwa API używa Django REST Framework do udostępniania funkcjonalności systemu przez RESTful API.

\subsubsection{ViewSets}

\begin{lstlisting}[language=Python, caption=RNGViewSet - główny endpoint dla RNG]
class RNGViewSet(viewsets.ViewSet):
    """ViewSet obsługujący operacje CRUD na generatorach RNG"""

    def __init__(self, **kwargs):
        super().__init__(**kwargs)
        # Dependency Injection - tworzenie zależności
        self.rng_repository = DjangoRNGRepository()
        self.result_repository = DjangoTestResultRepository()
        self.runners = [
            PythonRNGRunner(),
            ExeRNGRunner()
        ]

    def list(self, request):
        """GET /api/rngs - lista wszystkich generatorów"""
        rngs = self.rng_repository.get_all()
        serializer = RNGSerializer(rngs, many=True)
        return Response(serializer.data)

    def retrieve(self, request, pk=None):
        """GET /api/rngs/{id} - szczegóły generatora"""
        rng = self.rng_repository.get_by_id(int(pk))
        if not rng:
            return Response(
                {"error": "RNG not found"},
                status=status.HTTP_404_NOT_FOUND
            )
        serializer = RNGSerializer(rng)
        return Response(serializer.data)

    def create(self, request):
        """POST /api/rngs - utworzenie nowego generatora"""
        serializer = RNGSerializer(data=request.data)
        serializer.is_valid(raise_exception=True)

        rng = serializer.save_to_entity()
        saved_rng = self.rng_repository.save(rng)

        return Response(
            RNGSerializer(saved_rng).data,
            status=status.HTTP_201_CREATED
        )

    @action(detail=True, methods=['post'])
    def run_test(self, request, pk=None):
        """
        POST /api/rngs/{id}/run_test
        Uruchamia test statystyczny na generatorze
        """
        # Walidacja danych wejściowych
        test_name = request.data.get('test_name')
        samples_count = request.data.get('samples_count', 10000)
        seed = request.data.get('seed')
        test_parameters = request.data.get('test_parameters')

        # Wywołanie Use Case
        use_case = RunRNGTestUseCase(
            self.rng_repository,
            self.result_repository,
            self.runners
        )

        try:
            result = use_case.execute(
                rng_id=int(pk),
                test_name=test_name,
                samples_count=samples_count,
                seed=seed,
                test_parameters=test_parameters
            )

            serializer = TestResultSerializer(result)
            return Response(serializer.data)

        except Exception as e:
            return Response(
                {"error": str(e)},
                status=status.HTTP_500_INTERNAL_SERVER_ERROR
            )

    @action(detail=True, methods=['post'])
    def generate(self, request, pk=None):
        """
        POST /api/rngs/{id}/generate
        Generuje surowe bity bez przeprowadzania testów
        """
        count = request.data.get('count', 10000)
        seed = request.data.get('seed')

        rng = self.rng_repository.get_by_id(int(pk))
        if not rng:
            return Response(
                {"error": "RNG not found"},
                status=status.HTTP_404_NOT_FOUND
            )

        # Znajdź runner i wygeneruj
        runner = self._find_runner(rng)

        start_time = time.time()
        bits, data_type = runner.generate_raw(rng, count, seed)
        execution_time = (time.time() - start_time) * 1000

        return Response({
            "bits": bits,
            "count": len(bits),
            "execution_time_ms": execution_time,
            "rng_id": rng.id,
            "rng_name": rng.name,
            "seed": seed
        })
\end{lstlisting}

\subsubsection{Serializers}

Serializatory konwertują encje domenowe do/z formatu JSON:

\begin{lstlisting}[language=Python, caption=DRF Serializers]
class RNGSerializer(serializers.Serializer):
    """Serializer dla encji RNG"""

    id = serializers.IntegerField(read_only=True)
    name = serializers.CharField(max_length=200)
    language = serializers.ChoiceField(
        choices=[(lang.value, lang.name) for lang in Language]
    )
    algorithm = serializers.ChoiceField(
        choices=[(alg.value, alg.name) for alg in Algorithm]
    )
    description = serializers.CharField(allow_blank=True)
    code_path = serializers.CharField(max_length=500)
    parameters = serializers.JSONField(required=False, allow_null=True)
    is_active = serializers.BooleanField(default=True)
    created_at = serializers.DateTimeField(read_only=True)

    def save_to_entity(self) -> RNG:
        """Konwersja validated_data do encji domenowej"""
        return RNG(
            name=self.validated_data['name'],
            language=Language(self.validated_data['language']),
            algorithm=Algorithm(self.validated_data['algorithm']),
            description=self.validated_data.get('description', ''),
            code_path=self.validated_data['code_path'],
            parameters=self.validated_data.get('parameters'),
            is_active=self.validated_data.get('is_active', True)
        )

class TestResultSerializer(serializers.Serializer):
    """Serializer dla wyników testów"""

    id = serializers.IntegerField(read_only=True)
    rng_id = serializers.IntegerField()
    test_name = serializers.CharField()
    passed = serializers.BooleanField()
    score = serializers.FloatField(min_value=0.0, max_value=1.0)
    execution_time_ms = serializers.FloatField()
    samples_count = serializers.IntegerField()
    statistics = serializers.JSONField()
    seed = serializers.IntegerField(allow_null=True, required=False)
    test_parameters = serializers.JSONField(allow_null=True, required=False)
    created_at = serializers.DateTimeField(read_only=True)
\end{lstlisting}

\newpage

\subsection{Testy Statystyczne}

System implementuje kompletne baterie testów statystycznych dla weryfikacji jakości generatorów liczb pseudolosowych.

\subsubsection{NIST Statistical Test Suite}

Pełna implementacja 15 testów NIST SP 800-22:

\begin{table}[H]
\centering
\caption{Testy NIST - kompletna bateria}
\begin{tabular}{|l|p{6cm}|c|}
\hline
\textbf{Test} & \textbf{Opis} & \textbf{Min. bitów} \\
\hline
Monobit & Proporcja 0s i 1s & 100 \\
\hline
Block Frequency & Balans bitów w blokach & 100 \\
\hline
Runs & Liczba ciągów jedynek/zer & 100 \\
\hline
Longest Run & Najdłuższy ciąg jedynek & 128 \\
\hline
Binary Matrix Rank & Ranga macierzy binarnej & 38,912 \\
\hline
DFT (Spectral) & Analiza częstotliwości (FFT) & 1,000 \\
\hline
Non-Overlapping Template & Wzorce bez nakładania & 1,000 \\
\hline
Overlapping Template & Wzorce z nakładaniem & 1,000 \\
\hline
Universal (Maurer) & Test kompresji & 387,840 \\
\hline
Linear Complexity & Złożoność liniowa (Berlekamp-Massey) & 1,000,000 \\
\hline
Serial & Częstość m-bitowych wzorców & 1,000 \\
\hline
Approximate Entropy & Entropia przybliżona & 100 \\
\hline
Cumulative Sums & Sumy kumulacyjne & 100 \\
\hline
Random Excursions & Spacer losowy - cykle & 1,000,000 \\
\hline
Random Excursions Variant & Spacer losowy - warianty & 1,000,000 \\
\hline
\end{tabular}
\end{table}

\subsubsection{Diehard Battery of Tests}

Implementacja 15 klasycznych testów Diehard:

\begin{table}[H]
\centering
\caption{Testy Diehard - kompletna bateria}
\begin{tabular}{|l|p{7cm}|}
\hline
\textbf{Test} & \textbf{Opis} \\
\hline
Birthday Spacings & Kolizje dat urodzin w próbkach \\
\hline
Overlapping Permutations & Analiza nakładających się permutacji \\
\hline
Binary Rank (31x31) & Ranga macierzy binarnej 31x31 \\
\hline
Binary Rank (32x32) & Ranga macierzy binarnej 32x32 \\
\hline
Binary Rank (6x8) & Ranga macierzy binarnej 6x8 \\
\hline
Bitstream & Nakładające się 20-bitowe słowa \\
\hline
OPSO & Overlapping-Pairs-Sparse-Occupancy \\
\hline
OQSO & Overlapping-Quadruples-Sparse-Occupancy \\
\hline
DNA & Zliczanie 10-literowych "słów DNA" \\
\hline
Count the 1s (stream) & Liczenie jedynek w strumieniu \\
\hline
Count the 1s (byte) & Liczenie jedynek w bajtach \\
\hline
Parking Lot & Test parkowania - kolizje współrzędnych \\
\hline
Minimum Distance & Minimalna odległość między punktami \\
\hline
3D Spheres & Minimalna odległość punktów na sferze 3D \\
\hline
Squeeze & Test kompresji iteracyjnej \\
\hline
Overlapping Sums & Sumy nakładających się sekwencji \\
\hline
Runs & Analiza długości ciągów \\
\hline
Craps & Symulacja gry w kości \\
\hline
\end{tabular}
\end{table}

\subsubsection{Implementacja testów}

Wszystkie testy zwracają zunifikowany wynik:

\begin{lstlisting}[language=Python, caption=Struktura wyniku testu]
{
    "passed": bool,           # Czy test przeszedł (p-value > 0.01)
    "score": float,           # Wynik 0-1 (wyższy = lepszy)
    "statistics": {
        "p_value": float,     # P-value testu
        # ... inne statystyki specyficzne dla testu
    },
    "test_time_ms": float     # Czas wykonania testu
}
\end{lstlisting}

\newpage

\subsection{Konfiguracja Django}

\subsubsection{Settings.py}

Kluczowe ustawienia projektu:

\begin{lstlisting}[language=Python, caption=Konfiguracja Django]
# Backend Settings
INSTALLED_APPS = [
    'django.contrib.admin',
    'django.contrib.auth',
    'django.contrib.contenttypes',
    'django.contrib.sessions',
    'django.contrib.messages',
    'django.contrib.staticfiles',

    # Third-party apps
    'rest_framework',
    'corsheaders',

    # Local apps
    'io_rng',
]

MIDDLEWARE = [
    'corsheaders.middleware.CorsMiddleware',  # CORS - musi być na początku
    'django.middleware.security.SecurityMiddleware',
    'django.contrib.sessions.middleware.SessionMiddleware',
    'django.middleware.common.CommonMiddleware',
    'django.middleware.csrf.CsrfViewMiddleware',
    'django.contrib.auth.middleware.AuthenticationMiddleware',
    'django.contrib.messages.middleware.MessageMiddleware',
    'django.middleware.clickjacking.XFrameOptionsMiddleware',
]

# CORS Configuration - dla frontendu React
CORS_ALLOWED_ORIGINS = [
    "http://localhost:5173",      # Vite dev server
    "http://127.0.0.1:5173",
]

CORS_ALLOW_METHODS = [
    'DELETE',
    'GET',
    'OPTIONS',
    'PATCH',
    'POST',
    'PUT',
]

# REST Framework Configuration
REST_FRAMEWORK = {
    'DEFAULT_RENDERER_CLASSES': [
        'rest_framework.renderers.JSONRenderer',
    ],
    'DEFAULT_PARSER_CLASSES': [
        'rest_framework.parsers.JSONParser',
    ],
}

# Database - SQLite
DATABASES = {
    'default': {
        'ENGINE': 'django.db.backends.sqlite3',
        'NAME': BASE_DIR / 'db.sqlite3',
    }
}
\end{lstlisting}

\subsubsection{URL Routing}

\begin{lstlisting}[language=Python, caption=Konfiguracja URL]
# backend/io_rng/urls.py
from django.urls import path, include
from rest_framework.routers import DefaultRouter
from .api.views import RNGViewSet, TestResultViewSet

router = DefaultRouter()
router.register(r'rngs', RNGViewSet, basename='rng')
router.register(r'test-results', TestResultViewSet, basename='testresult')

urlpatterns = [
    path('api/', include(router.urls)),
]
\end{lstlisting}

\subsection{System Migracji Bazy Danych}

Django używa systemu migracji do zarządzania schematem bazy danych:

\begin{lstlisting}[language=bash, caption=Komendy migracji]
# Utworzenie nowych migracji po zmianie modeli
python manage.py makemigrations

# Zastosowanie migracji do bazy danych
python manage.py migrate

# Wyświetlenie stanu migracji
python manage.py showmigrations

# Cofnięcie migracji
python manage.py migrate io_rng 0004_previous_migration
\end{lstlisting}

Przykład migracji:

\begin{lstlisting}[language=Python, caption=Przykładowa migracja Django]
# io_rng/infrastructure/migrations/0005_testresultmodel_test_parameters.py
from django.db import migrations, models

class Migration(migrations.Migration):
    dependencies = [
        ('io_rng', '0004_previous_migration'),
    ]

    operations = [
        migrations.AddField(
            model_name='testresultmodel',
            name='test_parameters',
            field=models.JSONField(blank=True, null=True),
        ),
    ]
\end{lstlisting}

\subsection{Zalety architektury backendowej}

\begin{tcolorbox}[colback=green!5!white,colframe=green!75!black,title=Kluczowe zalety]
\begin{enumerate}
    \item \textbf{Clean Architecture} - pełna separacja warstw, niezależność od frameworka
    \item \textbf{Testowalność} - łatwe mockowanie dzięki Dependency Injection
    \item \textbf{Rozszerzalność} - łatwe dodawanie nowych runnerów, testów, algorytmów
    \item \textbf{Type Safety} - type hints w całym projekcie
    \item \textbf{Wielojęzyczność} - obsługa RNG w Python, C++, Rust, C\#
    \item \textbf{Uniwersalność} - Universal Adapter automatycznie dopasowuje się do różnych sygnatur
    \item \textbf{Wydajność} - optymalizacja poprzez kompresję Base64, numpy
    \item \textbf{Kompletność} - pełne baterie NIST (15 testów) i Diehard (15 testów)
\end{enumerate}
\end{tcolorbox}

\newpage

\section{Architektura Frontendowa}

\subsection{Struktura komponentów React}

Frontend używa architektury komponentowej z podziałem funkcjonalnym:

\begin{tcolorbox}[colback=blue!5!white,colframe=blue!75!black,title=Struktura komponentów]
\begin{itemize}
    \item \textbf{Layout:}
    \begin{itemize}
        \item \texttt{NavigationSidebar} - nawigacja boczna
        \item \texttt{Footer} - stopka
        \item \texttt{Header} - nagłówek
    \end{itemize}
    
    \item \textbf{Pages (Routes):}
    \begin{itemize}
        \item \texttt{Dashboard} - pulpit główny
        \item \texttt{Generator} - generator liczb
        \item \texttt{Tests} - testy statystyczne
        \item \texttt{Results} - wyniki testów
        \item \texttt{Wiki} - dokumentacja algorytmów
        \item \texttt{Settings} - ustawienia
    \end{itemize}
    
    \item \textbf{UI Components (shadcn/ui):}
    \begin{itemize}
        \item Button, Card, Dialog, Select
        \item Table, Tabs, Progress
        \item Chart (Recharts)
    \end{itemize}
    
    \item \textbf{Context Providers:}
    \begin{itemize}
        \item \texttt{ThemeProvider} - motyw (dark/light)
        \item \texttt{LanguageProvider} - internacjonalizacja
        \item \texttt{TestResultsProvider} - stan wyników testów
    \end{itemize}
\end{itemize}
\end{tcolorbox}

\subsection{Wzorce frontendowe}

\subsubsection{Context API dla globalnego stanu}

\begin{lstlisting}[language=Python, caption=Test Results Context]
export const TestResultsProvider = ({ children }) => {
  const [results, setResults] = useState<TestResult[]>([]);
  const [loading, setLoading] = useState(false);
  
  const fetchResults = async () => {
    setLoading(true);
    const data = await api.getTestResults();
    setResults(data);
    setLoading(false);
  };
  
  return (
    <TestResultsContext.Provider value={{ 
      results, loading, fetchResults 
    }}>
      {children}
    </TestResultsContext.Provider>
  );
};
\end{lstlisting}

\subsubsection{Custom Hooks}

System używa custom hooks dla logiki wielokrotnego użytku:

\begin{itemize}
    \item \texttt{useRNGList()} - pobieranie listy generatorów
    \item \texttt{useTestExecution()} - uruchamianie testów
    \item \texttt{useLanguage()} - zarządzanie językiem interfejsu
    \item \texttt{useTheme()} - zarządzanie motywem
\end{itemize}

\section{Komunikacja Backend-Frontend}

\subsection{REST API Endpoints}

\begin{table}[H]
\centering
\begin{tabular}{|l|l|p{6cm}|}
\hline
\textbf{Metoda} & \textbf{Endpoint} & \textbf{Opis} \\ \hline
GET & /api/rngs & Lista wszystkich RNG \\ \hline
POST & /api/rngs & Utworzenie nowego RNG \\ \hline
GET & /api/rngs/\{id\} & Szczegóły RNG \\ \hline
PUT & /api/rngs/\{id\} & Aktualizacja RNG \\ \hline
DELETE & /api/rngs/\{id\} & Usunięcie RNG \\ \hline
POST & /api/rngs/\{id\}/run\_test & Uruchomienie testu \\ \hline
POST & /api/rngs/\{id\}/generate & Generowanie liczb bez testu \\ \hline
GET & /api/rngs/\{id\}/test\_results & Wyniki testów dla RNG \\ \hline
GET & /api/test-results & Lista wszystkich wyników \\ \hline
GET & /api/test-results/\{id\} & Szczegóły wyniku \\ \hline
DELETE & /api/test-results/\{id\} & Usunięcie wyniku \\ \hline
\end{tabular}
\caption{REST API Endpoints}
\end{table}

\subsection{Przykładowe zapytanie i odpowiedź}

\textbf{Request:}
\begin{lstlisting}[language=bash]
POST /api/rngs/1/run_test
Content-Type: application/json

{
  "test_name": "chi_square",
  "samples_count": 100000,
  "seed": 42,
  "parameters": {
    "a": 1103515245,
    "c": 12345,
    "m": 2147483648
  }
}
\end{lstlisting}

\textbf{Response:}
\begin{lstlisting}[caption=JSON Response]
{
  "id": 123,
  "rng_id": 1,
  "test_name": "chi_square",
  "passed": true,
  "score": 0.87,
  "execution_time_ms": 234.56,
  "samples_count": 100000,
  "statistics": {
    "chi_square_statistic": 12.34,
    "p_value": 0.87,
    "degrees_of_freedom": 9
  },
  "created_at": "2026-01-14T10:30:00Z"
}
\end{lstlisting}

\section{Podsumowanie}

\subsection{Zastosowane wzorce architektoniczne}

\begin{enumerate}
    \item \textbf{Clean Architecture} - podział na warstwy: Core, Infrastructure, API
    \item \textbf{Repository Pattern} - abstrakcja dostępu do danych
    \item \textbf{Adapter Pattern} - Universal RNG Adapter
    \item \textbf{Strategy Pattern} - wymienne runnery dla języków
    \item \textbf{Use Case Pattern} - enkapsulacja logiki biznesowej
    \item \textbf{Dependency Injection} - wstrzykiwanie zależności
    \item \textbf{MVC/MVP} - separacja widoku i logiki (Frontend)
\end{enumerate}

\subsection{Zasady OOP}

\begin{enumerate}
    \item \textbf{SOLID:}
    \begin{itemize}
        \item Single Responsibility - każda klasa ma jedną odpowiedzialność
        \item Open/Closed - otwarte na rozszerzenia, zamknięte na modyfikacje
        \item Liskov Substitution - implementacje podstawialne
        \item Interface Segregation - minimalistyczne interfejsy
        \item Dependency Inversion - zależność od abstrakcji
    \end{itemize}
    
    \item \textbf{Enkapsulacja} - dane i logika ukryte w klasach
    \item \textbf{Polimorfizm} - różne implementacje tego samego interfejsu
    \item \textbf{Abstrakcja} - ukrycie szczegółów implementacji
    \item \textbf{Dziedziczenie} - przez interfejsy (ABC)
\end{enumerate}

\subsection{Struktury danych}

\begin{enumerate}
    \item \textbf{Dataclasses} - encje domenowe (RNG, TestResult)
    \item \textbf{Enumy} - typy wyliczeniowe (Language, Algorithm, DataType)
    \item \textbf{Dict} - parametry dynamiczne, JSON
    \item \textbf{List[T]} - kolekcje typowane
    \item \textbf{Tuple} - zwracanie wielu wartości
    \item \textbf{Optional[T]} - wartości nullable
    \item \textbf{Django Models} - ORM dla persystencji
    \item \textbf{JSONField} - elastyczne przechowywanie danych
\end{enumerate}

\subsection{Zalety architektury}

\begin{itemize}
    \item \textbf{Testowalność} - łatwe mockowanie interfejsów
    \item \textbf{Rozszerzalność} - łatwe dodawanie nowych runnerów, testów, algorytmów
    \item \textbf{Niezależność od frameworka} - logika biznesowa w Core
    \item \textbf{Separacja warstw} - clear separation of concerns
    \item \textbf{Type Safety} - type hints w Python, TypeScript w frontend
    \item \textbf{Maintainability} - czytelna struktura, SOLID
\end{itemize}

\end{document}